\documentclass[titlepage,a4paper]{jsarticle}
\usepackage{../../../sty/import}% 各種パッケージインポート
\usepackage{../../../sty/title}% タイトルページの変更
\usepackage{../../../sty/answergrid}
%% タイトルページの変数
% レポートタイトル
\title{心理学特論第14回目出席票}
% 提出日
\expdate{\today}
% 科目名
\subject{心理学特論}
% 分野
\class{情報経営システム工学分野}
% 学年
\grade{M1}
% 学籍番号
\mynumber{24336488}
% 記述者
\author{本間三暉}

\begin{document}
% titleページ作成
\maketitle
\section*{keyword\#1}
サブリミナル効果
\section*{keyword\#2}
バーナム効果
\section*{☆心理学特論の第14回テーマ「犯罪の心理学」について，簡潔にまとめてみましょう！}
\section{根拠のないデマやディープフェイクによるニュースなど，明らかにウソと思われるものが拡散してしまうのはどうしてかと考えてみて下さい．
悪意のあるウソが広まってしまう理由について，あなたの思うところを述べて下さい．}
根拠のないデマやディープフェイクが広がる理由は，情報を本当だと信じて拡散する人だけでなく，「面白い」「驚いた」「ネタとして共有したい」と考え，深く確認せずに「いいね」や共有をする人がいるからだと思う．
実際に，偽情報や誤情報は，驚きや感情を引き起こす内容ほど共有されやすいことが指摘されている．
また，投稿者や共有者には冗談だと分かっていても，投稿が元の文脈から離れて広がることで，後から見た人が事実だと誤解する可能性がある．
このように，悪意のない軽い反応であっても，結果としてウソの情報を広めることにつながるのだと思う．

\section{令和の日本で起こった犯罪について，印象的だったもの（未解決事件も含む）について，なぜ印象に残ったのか，理由も併せて挙げて下さい．}
令和の日本で印象に残っている犯罪は，いわゆる闇バイトによる強盗や強盗殺人事件である．
闇バイトでは，SNS上で高額報酬などを示して実行役を募集し，応募者から身分証明書や家族に関する情報を提出させることがある．
その後，犯罪だと気付いて断ろうとしても，本人や家族への危害をほのめかして脅され，犯行グループから抜け出せなくなる場合がある．
明らかに怪しい募集であっても，一度関わったことをきっかけに，強盗や殺人のような重大な犯罪にまで加担してしまう仕組みに恐ろしさを感じ印象に残った．
\end{document}
